% 第二章 V4：依据用户批注及第一章到第二章的阅读顺序重新组织。
% 保持七节、十二个三级小节；参考文献为本章讨论稿的独立编号。
\chapter{相关理论与技术基础}
\label{ch:foundation}

\section{本章引言}

端侧智能设备在处理图像、视频及其他感知信息时，需要根据应用的时间要求完成计算。不同推理应用具有各自的计算流程，同一设备又可能同时承担多个应用的请求。任务的执行时间不仅取决于计算内容，还与采用的处理器和工作频率有关；任务运行产生的功率则影响芯片温度，并进一步改变运行过程中的能量消耗。研究这类系统中的热感知任务调度，需要先了解计算平台、任务表示、能耗管理和温度变化的基本规律。

本章首先介绍实时系统和端侧异构多核处理器的基本特点，随后从经典实时任务出发，逐步介绍具有依赖关系的图任务以及深度学习推理的计算阶段和场景组织。在此基础上，分析处理器功耗及其能量管理技术，利用热阻热容网络描述多核系统的温度变化，最后介绍异构任务调度与优化求解的常用方法。这些内容为后续研究建立统一的理论基础。

\section{系统模型}

\subsection{实时系统模型}

实时系统不仅要求计算结果正确，还要求结果在规定时间内产生。例如，在工业控制中，系统需要在限定时间内根据传感器数据调整执行机构；在交通环境感知中，检测结果也需要及时送达后续决策环节。若计算完成时已经错过了必要的响应时机，即使结果本身准确，也可能无法满足应用要求。因此，任务的截止期限与计算结果的正确性共同构成实时系统的基本要求~[1]。

早期实时计算主要面向过程控制、航空航天及工业自动化等应用。随着移动计算和端侧人工智能的发展，实时性要求逐渐扩展到视频分析、人机交互和持续环境感知等场景。端侧设备通常在有限的计算资源与功耗条件下运行，摄像头产生的连续图像、传感器触发的检测任务以及用户发起的交互请求，都可能要求推理在一定时间内完成。当多个应用同时运行时，处理器资源的竞争还会使请求的响应时间发生变化。由此，实时调度不仅要安排任务的执行先后，还要考虑设备能力能否满足应用的时间需求~[2--4]。

不同应用对超时的容忍程度有所差异。根据截止期限违反后产生的后果，实时系统通常可分为硬实时、软实时和具有受限超时要求的系统~[1,18,19]。

（1）\textbf{硬实时系统。} 硬实时系统要求规定的每次任务执行都在截止期限之前完成，偶发超时也可能造成不可接受的后果。例如，若感知结果直接参与具有严格时限的设备保护或安全控制，系统就需要为相关计算提供明确的时间保证。此类系统通常依靠最坏情况执行时间估计及可调度性分析，验证所考虑运行条件下的期限要求。保证能够覆盖哪些干扰和执行变化，取决于分析所采用的任务与平台模型。

（2）\textbf{软实时系统。} 软实时系统同样关注响应时间，但少量任务未能按期完成时，主要表现为服务质量下降，而不一定导致整个系统失效。例如，视频内容分类或普通交互式图像处理偶尔出现响应延迟，可能影响画面的流畅性或用户体验。对于这类应用，平均响应时间、超时比例及整体服务质量通常也是重要评价指标。因此，软实时调度除了尽量满足各次请求的时限，还会关注系统在持续负载下的总体表现。

（3）\textbf{受限超时的实时系统。} 介于严格禁止超时和允许任意偶发超时之间，一些流式应用能够接受有限次数的结果丢失，却难以容忍连续多次未按期完成。例如，摄像头按固定帧率产生图像时，某一帧处理过晚，可以不再用于当前显示或跟踪；但若连续多帧未能及时处理，感知结果就会明显失去连续性。针对这类需求，$(m,k)$-firm约束要求任意连续$k$次作业中至少有$m$次按期完成；弱硬实时（Weakly-Hard Real-Time）模型则可进一步描述允许出现的超时模式~[18,19]。这些模型关注的不只是超时次数，还包括超时在时间上的分布。

上述分类反映了不同应用对任务完成时间的要求。对于需要逐次按期交付结果的推理服务，可以将每次请求的截止期限作为调度必须满足的条件；对于允许有限丢帧的服务，也可以采用相应的超时约束。本文研究的规划问题采用前一种要求：给定有限个请求及其执行参数，安排它们在各自截止期限内完成。具体任务如何表示、多个请求之间有哪些执行关系，将在下一节介绍。

\subsection{端侧异构多核处理器}

实时任务的时间要求需要依靠计算平台来满足。早期多核处理器主要将多个结构相近的CPU核心集成在同一芯片上，通过并行执行提高整体处理能力。这类同构多核平台适合使用相近的执行机制处理不同任务，但深度学习推理既包含控制与数据准备，又包含大量规则的矩阵及张量运算。单一类型处理器难以在这些计算环节中同时兼顾通用性、处理速度和能量效率，推动了通用核心与专用加速器协同工作的异构架构。

端侧片上系统（System on Chip，SoC）通常将中央处理单元（Central Processing Unit，CPU）、图形处理单元（Graphics Processing Unit，GPU）和神经网络处理单元（Neural Processing Unit，NPU）集成在同一平台，并通过片上互连与存储系统交换数据。CPU具有灵活的指令执行和分支控制能力，适合系统控制、数据准备及部分后处理工作；GPU由大量并行计算单元组成，适合数据并行程度较高的运算；NPU则针对神经网络常见的乘加计算配置专门的数据通路与片上存储，有利于高效执行部分卷积和矩阵计算~[2,3,7]。

不同处理器的计算能力互为补充，也使同一推理应用具有多种执行方式。例如，图像预处理可以由CPU完成，网络中的密集张量计算可以交由NPU执行，而某些不受NPU后端支持的计算仍可能需要GPU或CPU承担。实际部署应结合算子支持、执行时间、功率及数据交换情况选择资源，而不能简单地按照处理器种类决定所有计算阶段的执行位置。

图~\ref{fig:ch2_soc} 展示了一个由两个CPU执行单元、GPU和NPU组成的异构平台。不同计算单元通过片上互连交换数据，并且各自的资源占用和频率控制可能采用不同粒度。

\begin{figure}[htbp]
  \centering
  \begin{tikzpicture}[>=Latex,font=\small,
    unitdomain/.style={draw,rounded corners=4pt,thick,minimum height=2.2cm,align=center},
    block/.style={draw,rounded corners=2pt,fill=white,minimum width=1.18cm,minimum height=.65cm,align=center}]
    \draw[rounded corners=6pt,thick,fill=gray!4] (-.15,-1.68) rectangle (12.85,2.52);
    \node[font=\small\bfseries] at (6.35,2.16) {端侧异构片上系统（SoC）};
    \node[unitdomain,fill=blue!7,draw=blue!50!black,minimum width=3.70cm] (cpu) at (2.0,.7) {};
    \node[font=\bfseries] at (2,1.42) {CPU共享频率域};
    \node[block] at (1.1,.55) {CPU1};\node[block] at (2.9,.55) {CPU2};
    \node[font=\scriptsize] at (2,-.08) {分别执行任务};
    \node[unitdomain,fill=green!8,draw=green!45!black,minimum width=3.1cm] (gpu) at (6.32,.7) {};
    \node[font=\bfseries] at (6.32,1.42) {GPU频率域};
    \node[block,minimum width=2cm] at (6.32,.55) {GPU0};
    \node[font=\scriptsize] at (6.32,-.08) {并行运算};
    \node[unitdomain,fill=orange!8,draw=orange!50!black,minimum width=3.1cm] (npu) at (10.70,.7) {};
    \node[font=\bfseries] at (10.7,1.42) {NPU频率域};
    \node[block,minimum width=2cm] at (10.70,.55) {NPU0};
    \node[font=\scriptsize] at (10.7,-.08) {张量加速};
    \node[draw,rounded corners=2pt,fill=gray!18,minimum width=11.6cm,minimum height=.47cm] (bus) at (6.35,-1.23) {片上互连与存储资源};
    \foreach \x in {2.0,6.32,10.70}{\draw[-Latex] (\x,-.40) -- (\x,-.95);}
  \end{tikzpicture}
  \caption{端侧异构多核处理器的计算单元与频率控制域}
  \label{fig:ch2_soc}
\end{figure}

对调度而言，需要区分任务占用的计算资源和控制这些资源工作频率的范围。图中的CPU1与CPU2能够分别承担任务，但它们属于同一个频率控制域；两者并行工作时，需要使用相容的频率配置。相比之下，GPU和NPU可以分别调整工作点。因而，执行单元描述任务实际由哪里执行，频率控制域则描述哪些执行单元需要共同调节频率。

这种资源组织方式直接影响并发执行。映射到同一执行单元的任务通常需要先后使用资源；映射到不同单元的任务虽然可以并行，仍可能受到共享频率或数据依赖的限制。异构平台因此为协同推理提供了多种资源配置选择，同时也要求调度器合理协调处理器分配与执行时间。

\section{任务模型}

\subsection{实时任务模型}

系统能够在何时提供处理结果，取决于各项计算任务的到达时间和所需执行时间。实时调度理论通常使用任务释放时间、执行时间和截止期限描述这些要求。对于反复到达的任务，还需要考虑相邻两次计算请求之间的时间间隔。

经典实时任务模型以周期任务为基础。设周期任务$\tau_i$按固定时间间隔重复触发，可写为
\begin{equation}\label{eq:ch2-periodic}
  \tau_i=(C_i,T_i,D_i),
\end{equation}
其中，$C_i$表示任务的最坏情况执行时间，$T_i$表示相邻两次释放之间的周期，$D_i$表示相对截止期限。当$D_i=T_i$时，任务具有隐式截止期限。周期任务模型结构简单，便于分析长期负载，并为固定优先级和最早截止期限优先等实时调度方法提供基础~[1,5]。

周期任务描述一项反复执行的工作，而每一次触发都对应一次需要实际完成的计算。例如，摄像头每隔固定时间采集一帧图像，每帧对应一次新的处理需求。在实时调度中，这一次具体执行称为任务作业。若$\tau_i$的第$k$次作业在$r_{i,k}$时刻释放，则它的绝对截止时刻为
\begin{equation}\label{eq:ch2-jobdeadline}
  D_{i,k}^{\mathrm{abs}}=r_{i,k}+D_i.
\end{equation}
设该次作业的完成时刻为$F_{i,k}$，按期完成要求$F_{i,k}\le D_{i,k}^{\mathrm{abs}}$。释放时刻规定了计算最早可以开始的时间，绝对截止时刻则明确了这次计算需要完成的最晚时间。

在端侧应用中，计算请求并不总是严格周期到达。摄像头可以按帧率持续提供输入，也可能由传感器事件或用户操作触发额外的计算任务。对于后一类请求，给出每次请求的到达时刻和截止期限，比预先规定一个固定周期更直接。此外，当研究对象是有限时间范围内的一批已知请求时，可以逐一列出它们的释放时刻、执行需求及期限，并在这一范围内统一安排执行。这种表示既可容纳周期触发的请求，也适用于不规则的事件到达。

前面将一次请求看作一个完整任务，便于说明时间要求。但是，深度学习推理往往由多个计算步骤组成，各步骤还需要按照数据产生的顺序执行。为了利用多核处理器的并行能力，需要进一步描述一次任务内部的计算依赖关系。

\subsection{有依赖关系的图任务模型}

实际应用中的计算步骤往往具有先后关系。例如，图像数据需要先完成预处理，再进入不同的特征计算分支；当分支结果准备完毕后，系统才能进行融合并产生输出。若将整个过程视为不可拆分的任务，就无法描述哪些步骤必须等待，哪些步骤可以交由不同处理器并行完成。图结构任务模型由此成为描述复杂计算过程的一种常用方式~[6]。

图~\ref{fig:ch2_dag} 给出了一个具有分支与汇合关系的计算流程。任务A完成后，B与C可以分别处理不同部分的数据；任务D则需要同时获得B、C的结果后才能启动。图中的箭头表示计算的前后依赖，不代表两个任务必须使用同一执行单元。

\begin{figure}[htbp]
  \centering
  \begin{tikzpicture}[>=Latex,font=\small,
   node/.style={draw,rounded corners=3pt,minimum width=2.15cm,minimum height=.72cm,fill=blue!6,align=center},
   link/.style={->,thick}]
   \node[node] (a) at (0,0) {A：数据准备\\$w_A=3$};
   \node[node] (b) at (4.0,1.02) {B：分支计算\\$w_B=5$};
   \node[node] (c) at (4.0,-1.02) {C：分支计算\\$w_C=4$};
   \node[node] (d) at (8.0,0) {D：结果融合\\$w_D=2$};
   \draw[link] (a.east)--(b.west);\draw[link] (a.east)--(c.west);
   \draw[link] (b.east)--(d.west);\draw[link] (c.east)--(d.west);
   \draw[->,very thick,orange!80!black] (a.east)--(b.west);
   \draw[->,very thick,orange!80!black] (b.east)--(d.west);
   \node[font=\scriptsize,orange!70!black] at (4.0,1.73) {最长依赖路径：A--B--D};
  \end{tikzpicture}
  \caption{包含串行、并行分支及汇合关系的图任务示例（节点代价为示意值）}
  \label{fig:ch2_dag}
\end{figure}

当计算依赖不包含循环关系时，可采用有向无环图（Directed Acyclic Graph，DAG）描述任务。一个DAG记为
\begin{equation}\label{eq:ch2-dag}
  G=(V,E),\qquad E\subseteq V\times V,
\end{equation}
其中$V$是计算节点集合，$E$是节点之间的有向依赖边集合。对于$(i,j)\in E$，节点$j$需要等待节点$i$完成后才能开始。DAG不存在有向环，因此能够按照各节点的依赖关系形成合法的执行顺序。没有依赖关系的节点则可以在计算资源允许时并行执行。

除前后关系外，任务图中的每个节点还需要关联计算代价。设$w_i$表示节点$i$在选定的参考执行条件下所需的时间，则任务图的总计算代价为
\begin{equation}\label{eq:ch2-work}
  Q(G)=\sum_{i\in V}w_i.
\end{equation}
然而，总计算代价并不能单独决定整个应用的完成时间。图~\ref{fig:ch2_dag}中B和C能够并行，但D必须等待较晚结束的分支。如果四个节点的执行时间分别为3、5、4和2个时间单位，在不存在资源争用和额外通信时，全部工作量为14，而A--B--D依赖链的长度为10。即使有足够多的处理器，图任务也需要至少10个时间单位才能完成。

这一限制可以由最长依赖路径刻画。记$\mathcal P(G)$为从入口到出口的全部有向路径集合，则在节点代价固定时，关键路径长度为
\begin{equation}\label{eq:ch2-path}
  L(G)=\max_{\pi\in\mathcal P(G)}\sum_{i\in\pi}w_i.
\end{equation}
在暂不考虑通信开销且各节点采用给定执行时间的条件下，$L(G)$构成任务图完成时间的下界。对于异构平台，同一节点在不同处理器上的执行时间通常不同，因此关键路径的代价还会受到资源映射的影响。任务图保留的依赖与并行信息，为下一步面向具体推理应用进行阶段划分和处理器调度提供了基础。

\subsection{面向深度学习推理的任务模型}

深度学习推理的计算过程可以用图任务模型表示。神经网络由卷积、矩阵乘法、激活和归一化等运算构成，各运算通过张量数据传递形成先后关系。在原始计算图中，一个节点可能对应一个独立算子，也可能对应某种框架中的计算层。若直接以这些细粒度节点为单位安排异构处理器，调度器就需要分别处理大量算子，同时承担频繁的设备切换与中间数据传递。因此，推理引擎通常会根据算子特征和后端支持情况，通过算子融合、子图划分等方式组织实际执行过程~[3,7]。

例如，输入数据的格式转换与尺寸调整可以组成连续的前处理过程，卷积及其后的激活操作也可能在同一后端中融合执行。与其分别给这些基础操作安排处理器，不如把实际连续执行的部分作为较大的计算单位。这里将这样的单位称为计算阶段（Stage），并在不同计算阶段之间保留必要的数据依赖关系。计算阶段的划分由应用结构和部署方式确定，后续调度在该划分基础上选择处理器和执行时刻。

设神经网络原始算子集合为$V_{\mathrm{op}}$，将其划分为$K$个不重叠且非空的子集$S_k$，每个子集代表一个计算阶段，则有
\begin{equation}\label{eq:ch2-partition}
  \Pi=\{S_1,S_2,\ldots,S_K\},\qquad
  \bigcup_{k=1}^{K}S_k=V_{\mathrm{op}},\quad
  S_i\cap S_j=\varnothing\ (i\ne j).
\end{equation}
若某个子集中的算子需要另一子集的输出，就在对应计算阶段之间建立依赖边。由此得到的仍是一个DAG，只是节点粒度由原始算子变成了具有实际部署意义的计算阶段。

图~\ref{fig:ch2_stage}用一个简化的线性推理流程说明这一转换。左侧按运算步骤给出原始计算图；右侧则根据连续执行关系，将这些步骤组织为输入准备、主体计算和结果处理三个阶段。阶段数目及其划分方式可随推理应用和后端而变化，图中仅展示一种可能的组织方式。

\begin{figure}[htbp]
  \centering
  \begin{tikzpicture}[>=Latex,font=\small,
    op/.style={draw,rounded corners=2pt,fill=gray!5,minimum height=.52cm,minimum width=2.2cm,align=center},
    stage/.style={draw,rounded corners=3pt,thick,fill=blue!8,minimum height=.88cm,minimum width=2.6cm,align=center}]
    \node[font=\bfseries] at (1.5,2.75) {原始算子序列};
    \node[font=\bfseries] at (9.45,2.75) {按部署关系组织的计算阶段};
    \node[op] (o1) at (1.5,2.05) {尺寸调整};
    \node[op] (o2) at (1.5,1.34) {格式转换};
    \node[op] (o3) at (1.5,.63) {卷积运算};
    \node[op] (o4) at (1.5,-.08) {激活运算};
    \node[op] (o5) at (1.5,-.79) {分类映射};
    \node[op] (o6) at (1.5,-1.50) {结果筛选};
    \foreach \a/\b in {o1/o2,o2/o3,o3/o4,o4/o5,o5/o6}{\draw[->,gray!70] (\a.south)--(\b.north);}
    \draw[rounded corners=3pt,dashed,orange!60!black] (0.18,1.02) rectangle (2.82,2.41);
    \draw[rounded corners=3pt,dashed,orange!60!black] (0.18,-.42) rectangle (2.82,.98);
    \draw[rounded corners=3pt,dashed,orange!60!black] (0.18,-1.83) rectangle (2.82,-.46);
    \node[stage] (s1) at (9.45,1.65) {阶段1：输入准备};
    \node[stage] (s2) at (9.45,.05) {阶段2：主体计算};
    \node[stage] (s3) at (9.45,-1.50) {阶段3：结果处理};
    \draw[->,thick] (s1.south)--(s2.north);
    \draw[->,thick] (s2.south)--(s3.north);
    \draw[->,thick,orange!70!black] (2.95,1.72)--(7.8,1.72);
    \draw[->,thick,orange!70!black] (2.95,.24)--(7.8,.24);
    \draw[->,thick,orange!70!black] (2.95,-1.14)--(7.8,-1.14);
    \node[font=\scriptsize,fill=white] at (5.35,2.03) {阶段划分};
  \end{tikzpicture}
  \caption{由推理算子序列形成可调度计算阶段的示意过程}
  \label{fig:ch2_stage}
\end{figure}

推理应用的计算图描述了完成一次推理需要执行哪些阶段，而实际运行还需要考虑应用何时发起请求。对同一应用而言，每次输入到达都会产生一次新的推理请求。不同请求遵循相同的计算阶段结构，却需要分别执行各阶段并满足各自的截止期限。因此，需要区分应用中固定的计算阶段，以及某次请求中实际运行的阶段实例。前者决定计算关系，后者才占用运行期间的处理器资源。

当端侧设备同时运行多个应用时，不同应用的请求还可能在时间上重叠。例如，设备在一种使用状态下同时承担应用A与B的推理，在另一种状态下则运行A与C。这里把每种预先给定的应用组合及其请求集合称为一个场景。不同场景有各自的任务到达、资源竞争和时间要求，但同一应用的计算结构通常保持不变~[4]。

图~\ref{fig:ch2_scene}说明了两种场景的关系。应用A出现在两个场景中，示例中其主体计算阶段均部署在NPU0上；同时运行的B或C不同，因此A在两个场景中的实际开始时间、等待情况以及运行频率可以分别安排。应用A的固定部署在场景之间保持一致，而各场景的执行过程相互独立。

\begin{figure}[htbp]
  \centering
  \begin{tikzpicture}[font=\small,>=Latex,
    task/.style={draw,rounded corners=3pt,minimum width=2.4cm,minimum height=.70cm,align=center},
    shared/.style={draw,rounded corners=3pt,thick,fill=blue!11,minimum width=2.95cm,minimum height=.80cm,align=center}]
    \node[font=\bfseries] at (1.5,2.40) {应用组合};
    \node[font=\bfseries] at (5.85,2.40) {A的共同部署};
    \node[font=\bfseries] at (10.20,2.40) {场景内运行安排};
    \node[task,fill=green!7] (ab) at (1.5,1.27) {场景1：A+B};
    \node[task,fill=orange!8] (ac) at (1.5,-.65) {场景2：A+C};
    \node[shared] (a1) at (5.85,1.27) {A的主体计算阶段\\NPU0};
    \node[shared] (a2) at (5.85,-.65) {A的主体计算阶段\\NPU0};
    \node[task] (r1) at (10.20,1.27) {频率与时序方案1};
    \node[task] (r2) at (10.20,-.65) {频率与时序方案2};
    \draw[->,thick] (ab)--(a1);\draw[->,thick] (a1)--(r1);
    \draw[->,thick] (ac)--(a2);\draw[->,thick] (a2)--(r2);
    \draw[<->,thick,dashed,blue!65!black] (a1.south)--node[right,font=\scriptsize,fill=white] {同一具体处理器} (a2.north);
  \end{tikzpicture}
  \caption{多场景推理中的共同部署与独立运行配置示意}
  \label{fig:ch2_scene}
\end{figure}

这样，推理任务模型从单个请求的计算图扩展到多个请求和多种应用组合：任务图规定每次请求的计算与依赖，场景规定哪些请求共同使用设备，处理器部署及运行时序则决定这些计算如何落实到异构平台。进一步评价这些安排的能量开销，需要了解处理器功率随任务和工作状态变化的规律。

\section{多核系统能耗管理技术}

\subsection{功耗模型}

处理器运行需要消耗电能，且耗电速度会随着工作状态变化。对采用互补金属氧化物半导体（Complementary Metal-Oxide-Semiconductor，CMOS）工艺的数字电路，功耗一般分为动态功耗和静态功耗两部分。两者的产生机制不同，也会受到不同运行参数的影响。

动态功耗主要来自晶体管开关时的电容充放电。在常用近似下，其与电压、频率及开关活动程度的关系为
\begin{equation}\label{eq:ch2-dynamic}
  P_{\mathrm{dyn}}\approx\alpha_{\mathrm{sw}}C_{\mathrm{eff}}V^2f,
\end{equation}
其中$\alpha_{\mathrm{sw}}$为平均开关活动因子，$C_{\mathrm{eff}}$为等效开关电容，$V$为供电电压，$f$为时钟频率。对于同一处理器，当频率及其配套电压降低时，动态功耗通常随之下降。但不同计算任务会触发不同的电路活动和数据访问，同一频率下的实际功率也可能不同。

静态功耗主要来源于晶体管的漏电流。即使处理器暂时没有执行用户计算，部分器件仍可能存在漏电，维持时钟、缓存或其他系统功能也可能需要额外能量。漏电功耗可由
\begin{equation}\label{eq:ch2-leak}
  P_{\mathrm{leak}}=V I_{\mathrm{leak}}(V,T)
\end{equation}
表示，其中$I_{\mathrm{leak}}$与电压、器件特性以及工作温度$T$有关。动态功耗强调计算活动产生的开关消耗，静态功耗强调器件维持工作状态时的漏电消耗；两者共同构成处理器功耗的基本来源。

在任务调度中，研究者通常还需要区分处理器执行任务和处于空闲状态时的功率。执行任务会增加运算单元及存储系统的活动，空闲时则可能保留部分基础运行功率。因此，可以在规定的测量范围内，用空闲状态的基准功率和任务执行带来的额外功率描述负载变化。对于CPU、GPU和NPU，这些功率随工作点和计算内容而有所不同。如此描述关注的是调度时能够观察和计算的运行功率，与前面按电路产生机制划分的动态、静态功耗相互联系。

任务在时间轴上先后或并行运行时，系统功率$P(t)$也会相应变化。评价某一段运行过程的总能量，需要将该时间段内的功率累积起来，因此有
\begin{equation}\label{eq:ch2-energy}
  E(t_0,t_1)=\int_{t_0}^{t_1}P(t)\,\mathrm{d}t.
\end{equation}
若功率保持恒定，能量便等于功率与持续时间之积。对实际调度安排而言，处理器执行任务期间需要消耗能量，处于等待和空闲状态时也可能继续耗电。不同任务安排即使完成相同的计算，因运行频率、占用时间和空闲分布不同，系统能量仍可能不同。这也为通过运行状态管理降低能耗提供了可能。

\subsection{能耗管理及优化技术}

多核处理器的能耗管理主要利用运行负载和任务时间要求，在不同功率状态之间进行选择。常见方法包括动态电源管理（Dynamic Power Management，DPM）和动态电压频率调整（Dynamic Voltage and Frequency Scaling，DVFS）。DPM着重利用处理器暂时不承担计算的空闲间隔，DVFS则通过改变处理器的电压和频率协调计算速度与功率~[8,9]。

DPM在设备空闲时选择功耗较低的工作状态。例如，时钟门控可停止部分电路的开关活动，电源门控能够进一步降低被关闭电路的漏电。进入或退出这些状态可能需要额外时间和能量：如果空闲间隔很短，频繁切换反而未必能够节能。因此，DPM的状态选择与空闲时间长度、唤醒开销以及下一项任务的截止期限密切相关。

DVFS则通过选择处理器支持的电压—频率组合改变任务执行性能。提高频率通常能够缩短部分计算过程，却需要更高的活动功率；降低频率可以减少开关功耗，但也可能使任务占用设备更久。较低功率与较短执行时间分别有利于节能和及时释放资源，实际选择需要综合两者的变化。

不同任务的执行时间对频率变化并不具有完全相同的敏感程度。计算密集阶段通常更依赖处理器计算频率，而存在较多访存等待的阶段，其时长可能只在部分区间随频率明显改变。Da Costa在应用与硬件相关的性能模型中，将执行时间分为对频率敏感和相对不敏感的部分~[10]。若任务在参考频率$f^{\mathrm{ref}}$下的执行时长为$d^{\mathrm{ref}}$，这种关系可表示为
\begin{equation}\label{eq:ch2-response}
  d(f)=d^{\mathrm{ref}}\left[(1-\eta)+\eta\frac{f^{\mathrm{ref}}}{f}\right],
  \qquad 0\le\eta\le1,
\end{equation}
其中$\eta$表示在该模型中随频率反比变化的时长比例。当频率降为参考值的一半，任务时长为原来的$1+\eta$倍。若$\eta=1$，时长变成原来的两倍；若$\eta=0.3$，时长则为原来的1.3倍。即使两个阶段采用相同的降频幅度，也可能因计算特点不同而产生不同的时间代价。实际调频需从硬件允许的离散电压—频率工作点（Operating Performance Point，OPP）中选择。

图~\ref{fig:ch2_power_control} 从一个共同的时间窗口出发，示意DPM与DVFS如何改变任务执行和设备空闲期间的功率。上方安排较快完成工作，再利用较长的空闲间隔进入低功耗状态；下方则通过降低运行功率延长工作时间。图中阴影面积对应相应时段的能量，说明不同方法的结果取决于功率与持续时间的共同作用，而不是仅由运行时的功率高低决定。

\begin{figure}[htbp]
  \centering
  \begin{tikzpicture}[>=Latex,font=\small]
    \foreach \yy/\tag in {2.5/{方案一：较快运行 + 空闲管理},0/{方案二：低频运行}}{
      \draw[->] (2.2,\yy)--(12.35,\yy) node[right,font=\scriptsize] {时间};
      \draw[->] (2.2,\yy)--(2.2,\yy+1.72) node[above,font=\scriptsize] {功率};
      \node[align=right,font=\scriptsize] at (.98,\yy+.70) {\tag};
      \draw[densely dashed,gray] (11.35,\yy)--(11.35,\yy+1.55);
      \node[font=\scriptsize] at (11.35,\yy+1.78) {截止期限};
    }
    \fill[blue!17] (2.2,2.5) rectangle (5.45,3.82);
    \fill[gray!16] (5.45,2.5) rectangle (11.35,2.72);
    \draw[very thick,blue!65!black] (2.2,3.82)--(5.45,3.82)--(5.45,2.72)--(11.35,2.72);
    \node[font=\scriptsize] at (3.82,3.10) {执行任务};
    \node[font=\scriptsize] at (8.28,2.92) {空闲低功耗};
    \fill[green!17] (2.2,0) rectangle (9.85,.88);
    \fill[gray!16] (9.85,0) rectangle (11.35,.22);
    \draw[very thick,green!55!black] (2.2,.88)--(9.85,.88)--(9.85,.22)--(11.35,.22);
    \node[font=\scriptsize] at (6.0,.48) {执行时间延长};
    \node[font=\scriptsize] at (10.6,.42) {空闲};
  \end{tikzpicture}
  \caption{在共同时间窗口内利用空闲管理与调频改变功率分布的示意}
  \label{fig:ch2_power_control}
\end{figure}

在异构多核平台上，调频还受到频率控制域的影响。多个CPU执行单元并发运行时，可能需要使用同一工作频率；GPU和NPU则可以采用各自的合法工作点。某个阶段的降频会延长设备占用，进而影响其他阶段的开始时刻；较快完成则可能改变后续空闲时间。因此，合理的能耗管理既需要了解任务与频率之间的性能关系，也需要了解整个执行过程中各设备何时工作、何时空闲。

\section{多核系统热建模与温度分析}

\subsection{热阻热容网络模型}

处理器的运行功率不仅决定能量消耗，也影响芯片的温度。电路活动产生的热量会经芯片材料、封装和散热结构向外传递，同时还会在不同计算区域之间扩散。由于芯片及封装能够储存热量，处理器的功率发生变化后，温度通常不会立即达到新的平衡状态。这意味着当前温度既受此刻任务执行的影响，也与此前一段时间内的功率分布有关。

热阻热容（Resistance-Capacitance，RC）网络通过热学量与电路量之间的类比描述这一过程。在RC热模型中，温度差相当于电压差，热流相当于电流，热阻反映材料的导热阻力，热容反映相应区域储存热量的能力。Huang等提出的紧凑热模型将芯片与封装划分为相互连接的节点，为分析片上热点及瞬态温度提供了基础~[11]。

先考虑一个仅与环境交换热量的节点。设节点温度为$T(t)$，环境温度为$T_a$，热容为$C_{\mathrm{th}}$，向环境散热的热阻为$R_{\mathrm{th}}$，输入热功率为$p(t)$。根据节点储热量的变化等于注入热量减去散出热量，可得
\begin{equation}\label{eq:ch2-rc-one}
 C_{\mathrm{th}}\frac{\mathrm{d}T(t)}{\mathrm{d}t}
 =p(t)-\frac{T(t)-T_a}{R_{\mathrm{th}}}.
\end{equation}
等式右侧第一项是任务活动引起的发热，第二项是节点向环境传出的热量。当节点温度高于环境温度时，散热项随温差增加而增大；当输入功率与散热量相等时，节点进入热平衡状态。热容使温度变化具有一定的时间尺度，因此短时间的高功率运行与长时间的持续负载可能产生不同的温度轨迹。

单节点模型便于理解发热与散热过程，但多核SoC上的计算单元并不是彼此隔热的。例如，GPU持续执行高功率任务时，产生的热量可能传向相邻CPU和NPU区域，使尚未执行任务的单元也受到温度影响。因此，需要利用多个节点描述芯片内的热耦合。图~\ref{fig:ch2_rc} 展示了一个简化的四节点模型：CPU1、CPU2、GPU和NPU各具有局部热容，节点之间通过热阻连接，同时经各自散热路径向环境传递热量。

\begin{figure}[htbp]
 \centering
 \begin{tikzpicture}[>=Latex,font=\small,
    thermal/.style={draw,thick,circle,fill=orange!9,minimum size=1.42cm,align=center},
    source/.style={draw,rounded corners=2pt,fill=red!7,minimum width=1.48cm,minimum height=.55cm,align=center}]
   \node[thermal] (c1) at (0,1.20) {CPU1};
   \node[thermal] (c2) at (4.0,1.20) {CPU2};
   \node[thermal] (gpu) at (8.0,1.20) {GPU};
   \node[thermal] (npu) at (12.0,1.20) {NPU};
   \draw[<->,thick] (c1.east)--node[above,font=\scriptsize] {$R_{12}$}(c2.west);
   \draw[<->,thick] (c2.east)--node[above,font=\scriptsize] {$R_{2G}$}(gpu.west);
   \draw[<->,thick] (gpu.east)--node[above,font=\scriptsize] {$R_{GN}$}(npu.west);
   \node[source] (p1) at (0,-.57) {$p_1(t)$};
   \node[source] (p2) at (4,-.57) {$p_2(t)$};
   \node[source] (pg) at (8,-.57) {$p_G(t)$};
   \node[source] (pn) at (12,-.57) {$p_N(t)$};
   \foreach \a/\b in {p1/c1,p2/c2,pg/gpu,pn/npu}{\draw[->,thick] (\a.north)--(\b.south);}
   \node[draw,rounded corners=3pt,fill=blue!7,minimum width=12.2cm,minimum height=.55cm] (amb) at (6,3.55) {封装散热与环境温度 $T_a$};
   \foreach \a/\n in {c1/1,c2/2,gpu/G,npu/N}{\draw[->] (\a.north)--(\a.north|-amb.south) node[midway,right,font=\scriptsize] {$R_{\n a}$};}
   \node[font=\scriptsize] at (0,-1.11) {$C_1$};
   \node[font=\scriptsize] at (4,-1.11) {$C_2$};
   \node[font=\scriptsize] at (8,-1.11) {$C_G$};
   \node[font=\scriptsize] at (12,-1.11) {$C_N$};
 \end{tikzpicture}
 \caption{包含两个CPU温度节点及GPU、NPU节点的耦合RC热网络示意}
 \label{fig:ch2_rc}
\end{figure}

对一般的多节点网络，设节点$i$的热容为$C_i$，节点$i$与$j$之间的导热系数为$g_{ij}$，该节点与环境的导热系数为$g_{ia}$，其热平衡方程为
\begin{equation}\label{eq:ch2-rc-multi}
 C_i\dot T_i(t)=p_i(t)
 +\sum_{j\ne i}g_{ij}\bigl[T_j(t)-T_i(t)\bigr]
 +g_{ia}\bigl[T_a-T_i(t)\bigr].
\end{equation}
式中，节点间的温度差决定传热方向，而各节点的注热功率由对应区域的运行状态决定。进一步以$\boldsymbol\theta(t)=\boldsymbol T(t)-T_a\boldsymbol 1$表示节点温升向量，可将整个网络写为
\begin{equation}\label{eq:ch2-rc-matrix}
 \boldsymbol C\dot{\boldsymbol\theta}(t)
 =-\boldsymbol G\boldsymbol\theta(t)+\boldsymbol p(t),
\end{equation}
其中$\boldsymbol C$为热容矩阵，$\boldsymbol G$由节点间及对环境的导热关系构成，$\boldsymbol p(t)$为各节点接收的热功率。该表达同时保留了温度的时间变化和不同计算区域之间的传热作用。

\subsection{瞬态温度与热安全分析}

热阻热容模型描述的温度是一段连续变化的过程。为了直观理解这一特点，仍以单节点模型为例。如果节点在时刻$t_0$后的输入热功率$p$保持不变，则式~\eqref{eq:ch2-rc-one} 的解为
\begin{equation}\label{eq:ch2-single-response}
 \begin{aligned}
 T(t_0+\tau)=T_a+R_{\mathrm{th}}p
 +\bigl[T(t_0)-T_a-R_{\mathrm{th}}p\bigr]
 e^{-\tau/(R_{\mathrm{th}}C_{\mathrm{th}})}.
 \end{aligned}
\end{equation}
这里$R_{\mathrm{th}}C_{\mathrm{th}}$称为热时间常数。它说明温度从原有状态向新的平衡值变化需要一定时间，任务结束或功率下降后，温度也不会瞬间恢复到环境水平。一个高功率任务安排在较冷的初始状态下，与安排在前序任务已使芯片升温之后，其温度结果可能明显不同。

在多节点平台上，任务开始、结束及频率变化会使各区域注热功率发生变化，而节点间热传递又使温度变化相互影响。给定某段任务安排，可以先计算对应时间区间内的温度变化，再以该区间末端的温度作为下一区间的初始状态，沿整个任务执行过程传播。不同任务即使选用相同处理器，只要调整执行顺序或等待间隔，也可能得到不同的最高温度。F-DTM等研究已经利用未来任务信息和瞬态热预测协调CPU/GPU配置，体现了任务执行时序对热状态的重要影响~[14]。

温度还会通过器件漏电影响功耗。一般来说，处理器温度升高可能增加漏电流，使相同运行配置下的静态功率发生变化。在给定的有效温区内，可以围绕参考温度$T_i^{\mathrm{ref}}$，用线性近似描述某一节点相关漏电功率的变化：
\begin{equation}\label{eq:ch2-temp-leak}
 P_{i,\mathrm{leak}}(T_i)
 \approx P_{i,\mathrm{leak}}(T_i^{\mathrm{ref}})
 +\lambda_i\bigl(T_i-T_i^{\mathrm{ref}}\bigr),
\end{equation}
其中$\lambda_i$为相应状态与温度范围内的漏电敏感系数。功率影响温度，温度又影响漏电，由此形成相互关联的变化过程。Bao等与Zhou等的相关研究说明，温度变化及任务等待位置可以影响系统漏电能量与节能调度结果~[9,13]。

芯片温度还受到硬件安全上限的约束。设节点$i$允许的最高温度为$T_i^{\max}$，需要评价的时间区间为$[0,W]$，则热安全要求为
\begin{equation}\label{eq:ch2-thermal-safe}
 T_i(t)\le T_i^{\max},\qquad
 \forall t\in[0,W],\ \forall i.
\end{equation}
这一条件覆盖任务执行、设备等待和完成后的空闲阶段。对存在热耦合的多节点系统，最高温度可能出现在任务起止时刻之间，因此仅比较几个任务完成时的温度并不足以描述整个运行过程。Chantem等研究了实时多核任务的温度感知映射和调度，Schor等进一步针对实时应用分析运行期间可能达到的最坏温度，为将温度作为调度约束提供了理论依据~[12,15]。

至此，任务的处理器位置、开始时间和工作频率不仅决定计算何时完成，也决定各节点何时发热、如何散热以及是否满足温度要求。调度方法需要综合这些关系，在期限与热安全等条件下比较不同安排的能量与完成时间。

\section{异构任务调度与优化方法}

\subsection{异构最早完成时间算法}

任务图给出了计算依赖，异构平台则提供多个具有不同执行速度的处理器。任务调度需要综合这两方面信息，决定各节点按照什么顺序被安排，以及在哪个处理器上执行。列表调度是一类常见方法：先根据任务重要程度确定优先级，再按优先级依次为任务选择处理器。异构最早完成时间算法（Heterogeneous Earliest Finish Time，HEFT）是其中具有代表性的算法~[6]。

HEFT首先依据任务图的后继关系计算向上排名，用于衡量当前节点及其后续依赖路径上的剩余计算代价。设$\bar w_i$为节点$i$在候选处理器上的平均执行时间，$\bar c_{ij}$为节点$i$与$j$之间的平均通信代价，$\mathrm{succ}(i)$为$i$的直接后继集合，则
\begin{equation}\label{eq:ch2-heft-rank}
 \operatorname{rank}_u(i)=\bar w_i+
 \max_{j\in\mathrm{succ}(i)}
 \bigl\{\bar c_{ij}+\operatorname{rank}_u(j)\bigr\}.
\end{equation}
对于没有后继的出口节点，后面的最大项取零。向上排名较高的任务通常具有较长的后继计算路径，因此算法会优先处理这些节点。

接着，HEFT对当前准备安排的任务比较不同处理器上的最早完成时间。某个处理器即使执行该任务速度较快，如果之前安排的任务尚未结束，新的任务仍须等待。设$\operatorname{EST}(i,p)$表示任务$i$在处理器$p$上最早可以开始的时刻，$w_{ip}$表示对应的执行时间，则
\begin{equation}\label{eq:ch2-heft-eft}
 \operatorname{EFT}(i,p)
 =\operatorname{EST}(i,p)+w_{ip}.
\end{equation}
算法选择使EFT较小的处理器，并在该处理器已安排任务之间寻找合法空档。这里的空档插入是调度表的构造方式，已安排任务的执行不会因此被抢占。

例如，任务在CPU上只需2~ms，但CPU直到8~ms才空闲，则最早完成时间为10~ms；在GPU上需要3~ms，若GPU从4~ms起可用且前驱已完成，最早完成时间便是7~ms。HEFT会选择GPU，说明异构调度比较的是实际能够达到的完成时间，而不只是节点自身的执行时长。

HEFT通常能够较快构造一份利用处理器空档的时间表，适合作为调度初始化及性能导向的比较方法。它的经典选择准则主要关注完成时间；当任务还具有离散频率选项、系统能耗目标和温度限制时，则需要在候选评价中引入相应信息。这也引出了在更大配置空间中搜索调度方案的其他方法。

\subsection{元启发式优化算法}

列表调度通常按预定规则逐步形成一份安排，每一步选定的资源也会影响后续选择。对于包含较多设备和工作点组合的调度问题，还可以利用元启发式优化方法反复调整已有方案，以探索不同的映射与执行顺序。遗传算法（Genetic Algorithm，GA）和模拟退火（Simulated Annealing，SA）是常见代表。

遗传算法将任务映射、工作点等离散选择编码为一个个候选个体。首先生成初始种群，随后评价各候选的目标表现，通过选择、交叉和变异产生新的种群，并在多轮迭代中保留较有希望的方案。例如，同一个任务图可以对应若干套处理器配置，不同个体通过交换部分配置探索新的组合。当个体只记录离散选项时，还需通过调度构造过程确定实际启动时间，并检查任务依赖和资源占用情况。

模拟退火则从一个当前方案出发，反复生成少量改变后的邻域方案。若新方案较好便可以直接接纳；若稍差，也可能在搜索初期按一定概率接受，以增加探索范围。对于以能耗最小化为目标的问题，当候选能量增加$\Delta E>0$时，常用的接受概率为
\begin{equation}\label{eq:ch2-sa}
 P_{\mathrm{acc}}=\exp\!\left(-\frac{\Delta E}{\Theta}\right),
 \qquad \Delta E>0,
\end{equation}
其中$\Theta$是退火搜索的控制参数，随着迭代逐步降低。开始时较容易探索能量略高的方案，后期则逐渐偏向保留更好的候选。此处的退火温度与处理器的物理温度没有关系。模拟退火的基本思想来自Kirkpatrick等的组合优化研究~[16]。

元启发式方法可以灵活处理不同类型的配置变量，但搜索过程必须区分候选能量和候选是否符合约束。若某个安排因资源冲突、超出期限或超过温度上限而不可用，即使其计算能量较低，也不能直接作为可接受的调度结果。因而，搜索算法通常需要与完整时间表构造及相应的约束检查共同使用。

\subsection{基于数学规划的优化方法}

启发式方法通过一定规则探索候选，而数学规划将决策变量、目标和约束明确写成一个优化问题，再利用相应求解器进行计算。对于异构任务调度，处理器分配和离散频率选择属于离散决策，启动与完成时刻则属于时间决策。混合整数线性规划（Mixed-Integer Linear Programming，MILP）可以同时表示这两类变量，因此常用于具有设备选择、前驱关系和截止期限约束的调度问题~[8,17]。

例如，可以用二元变量记录某节点是否使用给定处理器和频点，并由这一选择关联执行时长；节点的启动时刻则受前驱完成及资源可用时间限制。类似地，同一处理器上的两个任务需要按照某种合法次序执行，每个请求的末端任务又需要满足相应截止期限。通过将这些要求表达成线性等式、不等式和整数变量条件，调度器便能够在共同目标下搜索满足约束的安排。

一般MILP可以表示为
\begin{equation}\label{eq:ch2-milp}
 \begin{aligned}
  \min_{\boldsymbol x}\quad &\boldsymbol c^{\mathsf T}\boldsymbol x,\\
  \text{s.t.}\quad &\boldsymbol A\boldsymbol x\le\boldsymbol b,\\
  &\boldsymbol A_{\mathrm{eq}}\boldsymbol x=\boldsymbol b_{\mathrm{eq}},\\
  &x_j\in\mathbb Z,\quad j\in\mathcal I,
 \end{aligned}
\end{equation}
其中$\mathcal I$为整数变量的下标集合。求解器可以利用找到的整数可行方案给出目标上界，同时通过松弛问题等途径形成下界。若上下界在相应精度要求下闭合，可以获得针对当前数学模型的最优性结论。对于变量组合较多的调度实例，找到可行方案与证明其最优可能需要不同的计算时间。

某些约束还可以在求解过程中逐步补充。约束生成方法通常先处理一个包含主要决策条件的模型，再独立检查当前候选是否满足其他要求；发现违规时，将能够排除相关违规安排的有效条件加入模型。以瞬态温度为例，一份满足设备占用和截止期限的时间表仍需要结合功率变化计算温度轨迹。若在某一节点发现明确超温，就可以利用该信息重新限制后续候选。具体能够构造何种有效约束，取决于所采用的热模型及调度变量表示。

此外，多场景推理中的部分配置需要在场景之间保持一致，其他任务则只影响所属场景。这种变量联系为分解优化提供了基础：先确定共同的处理器部署，再分别评价各场景中允许采用的频率和执行时序，并将局部结果合成为整体目标。这样的组织使跨场景共同决策与场景内任务调度具有清晰分工，也为后续设计热约束节能求解方法提供了思路。

\section{本章小结}

本章围绕端侧推理任务调度介绍了系统、任务、功耗和温度方面的相关理论。在系统模型中，首先从时间正确性出发介绍实时系统的主要类别，再说明CPU、GPU和NPU构成的异构平台如何提供不同类型的计算资源。在任务模型中，由周期任务及具体作业的时间要求出发，利用图模型描述计算依赖和关键路径，并进一步介绍推理计算阶段、请求实例及多场景应用组合。

随后，本章介绍动态功耗和静态功耗的基本规律、任务运行与空闲状态下的能量计算，以及DPM和DVFS的节能机制；通过单节点到多节点的RC网络说明芯片热耦合和瞬态温度的变化，并给出温度相关漏电及热安全条件。最后介绍HEFT、遗传算法、模拟退火和混合整数线性规划等调度与求解方法。上述理论为后续研究中协调处理器部署、运行配置、任务时序和连续温度约束提供了基础。

\section*{参考文献（第二章讨论稿引用对照）}
\addcontentsline{toc}{section}{参考文献（第二章讨论稿引用对照）}
\begingroup\small\setlength{\parindent}{0pt}\setlength{\parskip}{4pt}
[1] BUTTAZZO G C. Hard Real-Time Computing Systems: Predictable Scheduling Algorithms and Applications[M]. 3rd ed. Springer, 2011.\par
[2] KANG D, OH J, CHOI J, YI Y, HA S. Scheduling of Deep Learning Applications Onto Heterogeneous Processors in an Embedded Device[J]. IEEE Access, 2020, 8: 43980--43991.\par
[3] JEONG S, et al. Band: Coordinated Multi-DNN Inference on Heterogeneous Mobile Processors[C]. MobiSys, 2022.\par
[4] KIM J, HA S. Energy-Aware Scenario-Based Mapping of Deep Learning Applications Onto Heterogeneous Processors Under Real-Time Constraints[J]. IEEE Transactions on Computers, 2023, 72(6).\par
[5] LIU C L, LAYLAND J W. Scheduling Algorithms for Multiprogramming in a Hard-Real-Time Environment[J]. Journal of the ACM, 1973, 20(1): 46--61.\par
[6] TOPCUOGLU H, HARIRI S, WU M Y. Performance-Effective and Low-Complexity Task Scheduling for Heterogeneous Computing[J]. IEEE Transactions on Parallel and Distributed Systems, 2002, 13(3): 260--274.\par
[7] AGHAPOUR E, SAPRA D, PIMENTEL A D, PATHANIA A. PELSI: Power-Efficient Layer-Switched Inference[C]. RTCSA, 2023: 12--17.\par
[8] ESMAILI A, NAZEMI M, PEDRAM M. Modeling Processor Idle Times in MPSoC Platforms to Enable Integrated DPM, DVFS, and Task Scheduling Subject to a Hard Deadline[C]. ASP-DAC, 2019.\par
[9] BAO M, ANDREI A, ELES P, PENG Z. Temperature-Aware Idle Time Distribution for Leakage Energy Optimization[J]. IEEE Transactions on VLSI Systems, 2012, 20(7): 1187--1200.\par
[10] DA COSTA G. Hardware and Application Aware Performance, Power and Energy Models for Modern HPC Servers with DVFS[J]. Sustainable Computing: Informatics and Systems, 2025, 46: 101106.\par
[11] HUANG W, STAN M R, SKADRON K, et al. Compact Thermal Modeling for Temperature-Aware Design[C]. DAC, 2004: 878--883.\par
[12] CHANTEM T, HU X S, DICK R P. Temperature-Aware Scheduling and Assignment for Hard Real-Time Applications on MPSoCs[J]. IEEE Transactions on VLSI Systems, 2011, 19(10): 1884--1897.\par
[13] ZHOU J, et al. Thermal-Aware Correlated Two-Level Scheduling of Real-Time Tasks with Reduced Processor Energy on Heterogeneous MPSoCs[J]. Journal of Systems Architecture, 2018, 82: 1--11.\par
[14] MAITY S, ROY R, MAJUMDER A, et al. Future Aware Dynamic Thermal Management in CPU-GPU Embedded Platforms[C]. RTSS, 2022.\par
[15] SCHOR L, BACIVAROV I, YANG H, THIELE L. Worst-Case Temperature Guarantees for Real-Time Applications on Multi-Core Systems[C]. RTAS, 2012.\par
[16] KIRKPATRICK S, GELATT C D, VECCHI M P. Optimization by Simulated Annealing[J]. Science, 1983, 220(4598): 671--680.\par
[17] LIU W, YI J, LI M, CHEN P, YANG L. Energy-Efficient Application Mapping and Scheduling for Lifetime Guaranteed MPSoCs[J]. IEEE Transactions on Computer-Aided Design of Integrated Circuits and Systems, 2019.\par
[18] BERNAT G, BURNS A, LLAMOS\'{I} A. Weakly-Hard Real-Time Systems[J]. IEEE Transactions on Computers, 2001, 50(4): 308--321.\par
[19] HAMDAOUI M, RAMANATHAN P. A Dynamic Priority Assignment Technique for Streams with $(m,k)$-Firm Deadlines[J]. IEEE Transactions on Computers, 1995, 44(12): 1443--1451.\par
\endgroup
